%----------------------------------------------------------------------------------------
% Ilyes Jaouedi — CV, ADAS / Perception Engineer version
%----------------------------------------------------------------------------------------
%----------------------------------------------------------------------------------------
% Ilyes Jaouedi — shared CV base (preamble, helpers, common blocks)
% Included by cv-research.tex, cv-radar.tex, cv-adas.tex, cv-mle.tex
%----------------------------------------------------------------------------------------
\documentclass[a4paper, 10pt]{article}

\usepackage[T1]{fontenc}
\usepackage[utf8]{inputenc}
\usepackage[a4paper, top=0.9cm, bottom=0.9cm, left=1.7cm, right=1.7cm]{geometry}
\usepackage{mathpazo}
\linespread{1.02}
\usepackage{microtype}
\usepackage[hidelinks]{hyperref}
\usepackage{fontawesome5}
\usepackage{enumitem}
\usepackage{xcolor}

\pagestyle{empty}
\setlength{\parindent}{0pt}
\setlength{\parskip}{0pt}

\hypersetup{colorlinks=true, urlcolor=black}

%----------------------------------------------------------------------------------------
% HELPERS
%----------------------------------------------------------------------------------------

% Section header: text + rule filling to the right
\newcommand{\cvsection}[1]{%
  \vspace{5pt}%
  \par\noindent{\large #1}\hspace{0.5em}\hrulefill\par%
  \vspace{2pt}%
}

% Entry: #1 role  #2 organisation  #3 period  #4 location
\newcommand{\cventry}[4]{%
  \par\noindent\textbf{#1}\hfill\textit{#2}\\[1pt]%
  \textit{#3}\hfill{\small #4}\par%
  \vspace{1pt}%
}

% Publication row: fixed-width bold venue | authors | italic title
\newcommand{\cvpub}[3]{%
  \par\noindent\hangindent=10.8em\hangafter=1%
  \makebox[10.4em][l]{\textbf{#1}}\ifx&#2&\else#2.\enspace\fi\textit{#3}\par%
  \vspace{2pt}%
}

% Project row: #1 title  #2 short context (right-aligned)  #3 description
\newcommand{\cvproject}[3]{%
  \par\noindent\textbf{#1}\hfill{\small\textit{#2}}\\[1pt]%
  #3\par%
  \vspace{3pt}%
}

% Tight bullet list
\newenvironment{cvitems}{%
  \begin{itemize}[leftmargin=1.2em, itemsep=1pt, topsep=2pt, parsep=0pt]
}{%
  \end{itemize}\vspace{3pt}%
}

%----------------------------------------------------------------------------------------
% HEADER — #1 is the per-version tagline
%----------------------------------------------------------------------------------------
\newcommand{\cvheader}[1]{%
  \begin{center}
    {\huge Ilyes Jaouedi}\\[3pt]
    {\small\scshape #1}\\[5pt]
    {\small
      \faEnvelope\enspace\href{mailto:ilyes.jaouedi@centralesupelec.fr}{ilyes.jaouedi@centralesupelec.fr}
      \quad\textbar\quad
      \faGlobe\enspace\href{https://ilyes-jaouedi.github.io}{ilyes-jaouedi.github.io}
      \quad\textbar\quad
      \faGithub\enspace\href{https://github.com/ilyes-jaouedi}{ilyes-jaouedi}
      \quad\textbar\quad
      \faLinkedinIn\enspace\href{https://linkedin.com/in/ilyes-jaouedi}{ilyes-jaouedi}
    }\\[2pt]
    {\small\faMapMarker*\enspace Paris area, France
      \quad\textbar\quad
      \textbf{Available for a permanent position (CDI) from mid-2027}}
  \end{center}
  \vspace{1pt}%
}

%----------------------------------------------------------------------------------------
% SHARED BLOCKS
%----------------------------------------------------------------------------------------

% Compact education block: one line per degree
\newcommand{\cvedu}[3]{%
  \par\noindent\textbf{#1}, #2\hfill\textit{#3}\par\vspace{1pt}%
}
\newcommand{\EducationBlock}{%
  \cvsection{Education}
  \cvedu{PhD, Signal Processing \& AI}{Université Paris-Saclay (L2S, CentraleSupélec)}{2024 -- 2027}
  {\small\hspace{1em}\emph{Sensor Fusion via Optimal Transport for Driver Assistance}.
   Supervisors: G.~Chardon, J.~Picheral.}\par\vspace{2pt}
  \cvedu{Engineering Degree, Signal Processing \& Systems}{Polytechnique de Tunisie (ranked 1\textsuperscript{st}, final year)}{2021 -- 2024}
  \cvedu{Preparatory Cycle, Maths \& Physics}{IPEIM Monastir (ranked 38\textsuperscript{th}/1800 nationally)}{2019 -- 2021}
}

% Publications, patents and preprints
\newcommand{\PublicationsBlock}{%
  \cvsection{Publications \& Patents}
  \cvpub{IEEE ICASSP 2026}%
        {\textbf{I.~Jaouedi}, G.~Chardon, J.~Picheral}%
        {Sensor Array and Camera Fusion via Unbalanced Optimal Transport for 3D Source
         Localization. Barcelona. \href{https://hal.science/hal-05572212v1}{hal-05572212}}
  \cvpub{GRETSI 2025}%
        {\textbf{I.~Jaouedi}, G.~Chardon, J.~Picheral}%
        {Fusion Acoustique--Visuelle par Transport Optimal. Strasbourg.
         \href{https://hal.science/hal-05273990}{hal-05273990}}
  \cvpub{Patent (FR) 2026}%
        {}%
        {Two French patent applications with FORVIA and CentraleSupélec (fast 3D radar imaging;
         camera--radar fusion). Filing in progress.}
}
% Full version adds work in preparation (research CV)
\newcommand{\PublicationsBlockFull}{%
  \PublicationsBlock
  \cvpub{Under review}%
        {}%
        {Computer-vision paper on learning dense radar reflectivity from images.}
  \cvpub{In preparation}%
        {}%
        {IEEE journal extension of the ICASSP paper; two further papers pending patent filing.}
}

\newcommand{\CertificationBlock}{%
  \cvsection{Certifications}
  \cvpub{Google Cloud 2025}%
        {}%
        {Professional Machine Learning Engineer ---
         \href{https://www.credly.com/badges/55072af5-5964-455b-9310-0a770916e7f5/public_url}{Verified on Credly}}
  \cvpub{Google Cloud 2026}%
        {}%
        {Generative AI Leader ---
         \href{https://www.credly.com/badges/27745c52-ea91-48d1-8070-5bc5424ce931/public_url}{Verified on Credly};
         Associate Data Practitioner ---
         \href{https://www.credly.com/badges/f98a92be-76c2-45b0-b820-17678cf4b580/public_url}{Verified on Credly}}
}


\begin{document}

\cvheader{Perception Engineer --- Camera, Radar \& LiDAR Fusion for ADAS}

\noindent
PhD candidate in camera--radar fusion for ADAS (with FORVIA), with prior internships in
LiDAR--camera 3D detection and tracking. I have worked with camera, imaging radar and LiDAR,
from raw sensor data to fused 3D output. Two patent applications with FORVIA.
\vspace{1pt}

%----------------------------------------------------------------------------------------
\cvsection{Experience}

\cventry{PhD Candidate --- Camera--Radar Fusion for ADAS}%
        {L2S, CentraleSupélec / Univ.\ Paris-Saclay \,$\cdot$\, FORVIA}%
        {October 2024 -- Present}%
        {Gif-sur-Yvette, France}
\begin{cvitems}
  \item \textbf{Camera--radar fusion on RADIal:} built a training-free pipeline that fuses raw
        radar antenna signals with camera detections and segmentations (SAM). Point precision
        rises from 4.1\% (CFAR baseline) to 67.2\% at 85.7\% object recall, measured with a
        point-in-box protocol I designed. \emph{(Patent application.)}
  \item \textbf{High-resolution 3D radar imaging:} dense range--azimuth--elevation point clouds
        from raw ADC data, more than 10$\times$ faster than MUSIC-FBSS; validated on a TI IWR6843
        + DCA1000 in the FORVIA anechoic chamber ($<$4\,cm error) and on the 192-element RADIal
        array. \emph{(Patent application.)}
  \item \textbf{Radar-aware camera model:} predicts radar reflectivity from a single camera image
        (frozen DINOv3 + DPT decoder), trained on radar detections and undetected cells with a
        censored likelihood that removes the range bias; AUC 0.85 vs 0.73 for detection
        prediction. \emph{(Paper under review.)}
  \item \textbf{Sensor fusion with guarantees:} array--camera fusion via unbalanced optimal
        transport, with Cramér--Rao bounds proving the gain over each sensor alone.
        \emph{(IEEE ICASSP 2026.)}
\end{cvitems}

\cventry{Computer Vision Research Intern}%
        {CAOR, Mines Paris -- PSL}%
        {February 2024 -- July 2024}%
        {Paris, France}
\begin{cvitems}
  \item Reconstructed vehicle trajectories by fusing traffic-camera video with LiDAR; built
        3D reconstruction and tracking pipelines for urban traffic scenes.
\end{cvitems}

\cventry{Computer Vision Engineer Intern}%
        {Enova Robotics}%
        {Summer 2023}%
        {Sousse, Tunisia}
\begin{cvitems}
  \item Implemented, benchmarked and fine-tuned LiDAR--camera fusion models (PV-RCNN, BiProDet)
        for 3D object detection on custom robotic datasets.
\end{cvitems}

%----------------------------------------------------------------------------------------
\cvsection{Selected Projects}

\cvproject{Real-time Object Detection --- YOLOv8}{Deployment}%
  {Fine-tuning pipeline with TensorRT export and real-time video inference.}
\cvproject{RadarFlows --- Synthetic Radar Generation}{Data augmentation}%
  {Generates realistic radar maps from camera images and 3D detections, for data augmentation.}

%----------------------------------------------------------------------------------------
\cvsection{Skills}

\noindent
\textbf{Perception:} camera--radar and LiDAR--camera fusion, 3D object detection, object
  tracking, point clouds, segmentation, camera calibration and projection models\\
\textbf{Radar:} FMCW MIMO (TDM, DDM), range--Doppler--angle processing, CFAR, raw ADC data,
  TI mmWave, RCS\\
\textbf{Data:} RADIal (raw automotive radar + camera), own anechoic-chamber radar recordings\\
\textbf{ML \& deployment:} PyTorch, vision foundation models (SAM, DINOv2, Depth Anything),
  ONNX/TensorRT, OpenCV\\
\textbf{Programming:} Python, C++, MATLAB, Docker, Git, SLURM, Linux\\
\textbf{Languages:} Arabic (native), French (fluent), English (fluent)

\PublicationsBlock
\EducationBlock

\end{document}
